\documentclass[11pt]{article}

\usepackage[T1]{fontenc}
\usepackage[utf8]{inputenc}
\usepackage{lmodern}
\usepackage{geometry}
\geometry{margin=1in}
\usepackage{hyperref}
\usepackage{booktabs}
\usepackage{longtable}
\usepackage{array}
\usepackage{amsmath}
\usepackage{amssymb}
\usepackage{listings}
\usepackage{xcolor}
\setlength{\emergencystretch}{3em}

\hypersetup{
  colorlinks=true,
  linkcolor=blue!50!black,
  citecolor=blue!50!black,
  urlcolor=blue!50!black
}

\lstset{
  basicstyle=\ttfamily\small,
  breaklines=true,
  frame=single,
  columns=fullflexible
}

\newcommand{\code}[1]{\texttt{#1}}
\newcommand{\HE}{\textup{HE}}
\newcommand{\CeTTa}{\textup{CeTTa}}
\newcommand{\PathMap}{\textup{PathMap}}

\title{Faithful Backends for \HE{} MeTTa\\
\large Bag Semantics, Candidate Selection, Co-Reference Preservation,\\
and Revision-Based Query Caching}
\author{Zar Goertzel\\Oru\v{z}i}
\date{March 2026}

\begin{document}
\maketitle

\begin{abstract}
This note sketches a correctness story for backend implementations of
\HE{}~MeTTa.  The central claim is that a practical backend should not be
identified with the full runtime semantics of a MeTTa space.  The
semantics lives at the bag level: repeated derivations and repeated atoms
matter.  The backend may quotient that semantics to support-level
structure for indexing, as long as the quotient and the refinement layers
are explicit and proved faithful.
%
Our current Lean development decomposes this story into three layers:
lists of evaluator outcomes, multisets of semantic results, and finite
support/index views.  On top of this, we add a concrete candidate-selector
architecture for \PathMap{} tries, a counted-space refinement explaining
why duplicate adds need not change trie support, an injective
canonicalization discipline for variable identities, a variant-key query
correctness layer for alpha-equivalent cache keys, and a revision-based
cache model for query reuse.  Together these results explain why
\CeTTa{}'s runtime should use \PathMap{} for candidate narrowing,
native matching for exact unification, explicit count metadata for bag
semantics, variant-aware materialization for cached answers, and explicit
invalidation rules for cached query answers.
\end{abstract}

\section{Problem Statement}

Backend bugs in MeTTa runtimes tend to fall into a small number of
families:
\begin{enumerate}
\item candidate selectors that are too weak or too strong;
\item storage layers that silently collapse semantic multiplicity;
\item canonicalization passes that preserve alpha-shape but break
  co-reference;
\item query caches that reuse answers beyond the mutation boundary for
  which they were computed.
\end{enumerate}

These failures are easy to confuse because they often show up through the
same surface symptom: a query returns the wrong answer after spaces have
been indexed, cloned, canonicalized, or cached.  The right response is
not to ``prove the runtime'' in one monolithic theorem.  Instead we prove
one obligation per runtime seam.

\section{Semantic Layers}

The first part of the development says that the semantic carrier for
\HE{} evaluation is not a set but a bag.

\paragraph{Evaluator layer.}
The runtime naturally produces a list of outcomes.  This is the
computational layer where order, scheduling, and implementation strategy
still matter.

\paragraph{Semantic layer.}
\code{NondeterminismCarrier.lean} moves to \code{Multiset}, yielding the
correct semantic equality for repeated results.  \code{BagSupportBridge}
then proves that support is the quotient visible to set-like indexing.

\paragraph{Index layer.}
\code{MonadMorphismChain.lean} and \code{SpaceQuerySupport.lean} package
the resulting bridge:
\[
\texttt{List} \to \texttt{Multiset} \to \texttt{Finset}.
\]
This is the level at which a trie, hash table, or cached query store is
allowed to operate.

\paragraph{Positive example.}
Adding the same atom twice should increase the bag count while leaving the
support unchanged.

\paragraph{Negative example.}
Treating a set-valued index as if it were itself the full MeTTa space
semantics.

\section{Faithful Backend Contract}

\code{SpaceQuerySupport.lean} gives the contract point:
\begin{quote}
\code{FaithfulBackend} means backend queries agree with \HE{} queries at
support level.
\end{quote}

This is the right initial obligation for a backend such as a
\PathMap{}-backed space:
\begin{enumerate}
\item answer the same support-level equation queries;
\item answer the same support-level type queries;
\item refine support semantics with explicit count or row metadata, rather
  than burying multiplicity in trie keys.
\end{enumerate}

The counted-space refinement in
\code{CandidateArchitecture.lean} makes this concrete.  \PathMap{} stores
support.  The multiplicity layer stores how many copies exist.  The two
are intentionally separated.

\section{Candidate Selection and \PathMap{}}

Theorems in \code{CandidateArchitecture.lean} now justify the architecture
that \CeTTa{} uses in practice:
\begin{enumerate}
\item the trie performs candidate selection;
\item native matching filters those candidates exactly;
\item the composition equals direct querying.
\end{enumerate}

The abstract theorem \code{twoPhase\_eq\_direct} proves the shape of the
argument.  The concrete theorem \code{concrete\_twoPhase\_eq\_direct}
instantiates it for a trie-based selector.  The prefix narrowing results
then move beyond the unrealistic ``scan all support'' model and state the
performance-relevant theorem: useful selectors may be strict support
subsets without losing correctness, provided they are sound and
phantom-free.

This is the strongest answer to the question ``why not let the backend do
all the unification?''  Because candidate narrowing and exact matching are
different obligations.  \PathMap{} is excellent at the first and should
not be forced to simulate the second through a lossy bridge layer.

\section{Co-Reference Preservation}

The nastiest runtime bugs are usually not about support.  They are about
binders.

\code{CoReferencePreservation.lean} isolates the right invariant:
variable canonicalization must be injective.  Under an injective renaming:
\begin{enumerate}
\item repeated source variables stay repeated;
\item distinct source variables stay distinct;
\item composed renamings remain injective.
\end{enumerate}

This is the formal obligation behind runtime components such as
\code{term\_universe.c} and the canonicalization logic in
\code{table\_store.c}.  If a runtime stores
\[
(\to (\to \varphi \psi) (\to \varphi \varphi))
\]
and later materializes
\[
(\to (\to \varphi \psi) (\to \chi \chi)),
\]
then the bug is not ``just alpha-equivalence.''  It is a failure to
preserve co-reference.

Recent work strengthens this seam in two directions.  First,
\code{BindingComposition.lean} exposes \code{Bindings.Extends} as the
right invariant for seed-preserving matching.  Second,
\code{VariantQueryCorrectness.lean} proves a bisimulation theorem
\texttt{simpleMatch\_rename\_bisim} using the fuel-free renaming
\texttt{applyAtomTotal}.  This proves that alpha-equivalent query keys do
not merely preserve support; they preserve the right notion of matching
success needed for variant-key cache reuse.

\section{Revision-Based Query Caching}

\code{CacheCorrectness.lean} already captures the main cache seam:
\begin{enumerate}
\item populate a cache entry from a space and a query;
\item reuse it only when the space revision still matches;
\item invalidate lazily on mutation;
\item preserve cache validity for non-crossing mutations such as
  duplicate \code{add-atom-nodup}.
\end{enumerate}

The important point is that cache correctness is not merely a performance
claim.  It is part of backend faithfulness.  A cache that returns stale
support answers is semantically wrong even if the trie and matcher are
otherwise correct.

The newer files \code{VariantQueryCorrectness.lean} and
\code{BindingCacheSoundness.lean} extend this from support-level cache
discipline to binding-aware reuse:
\begin{enumerate}
\item \texttt{variant\_queries\_same\_rhs} proves that
  alpha-equivalent queries produce the same right-hand-side atoms after
  \code{queryEquations};
\item \texttt{variant\_cache\_binding\_shape} proves that the result
  shapes agree under variant reuse;
\item \texttt{cache\_binding\_exact} and related lemmas isolate the
  exact-versus-variant parts of the cache story.
\end{enumerate}

The theorem family around mutation invalidation is especially important
for implementations that mix:
\begin{enumerate}
\item candidate narrowing via a support index;
\item exact matching via native matcher state;
\item explicit count metadata;
\item snapshot or persistent clones.
\end{enumerate}

\section{Runtime Seam Map}

\begin{center}
\begin{tabular}{@{}p{0.28\linewidth}p{0.28\linewidth}p{0.34\linewidth}@{}}
\toprule
\textbf{Lean file} & \textbf{Runtime seam} & \textbf{Obligation} \\
\midrule
\code{CandidateArchitecture.lean} & \code{space\_match\_backend.c} & trie candidate selection + native exact rematch \\
\code{BagSupportBridge.lean} & counted-space layer & support/count split over bag semantics \\
\code{CoReferencePreservation.lean} & \code{term\_universe.c} & canonical storage preserves variable identity \\
\code{VariantQueryCorrectness.lean} & \code{table\_store.c} & variant-key queries preserve RHS answers \\
\code{BindingCacheSoundness.lean} & \code{table\_store.c} & cached answer shapes remain binding-compatible \\
\code{BindingComposition.lean} & \code{space\_match\_backend.c} & seed bindings survive rematch/merge \\
\code{SeedRematchContract.lean} & \code{space\_subst\_match\_with\_seed} & seeded rematch extends seed, rejects loops, imported binding parity \\
\code{CacheCorrectness.lean} & \code{table\_store.c} & revision-based query reuse and invalidation \\
\code{IncrementalTableSemantics.lean} & future \code{table\_store.c} & partial table soundness, answer monotonicity, completion \\
\code{CeTTaRuntimeContracts.lean} & \code{term\_canon.c}, \code{session.c} & canonicalization safety, rollback, effect/profile contracts \\
\code{SpaceQuerySupport.lean} & backend interface & support-level agreement target \\
\bottomrule
\end{tabular}
\end{center}

\section{What Belongs Here, and What Does Not}

This paper should carry the full backend-correctness story:
\begin{enumerate}
\item bag semantics and support quotients;
\item counted-space refinement;
\item candidate-selection correctness;
\item \textbf{seeded native rematch} --- the theorem that turns candidate parity into binding parity (\code{SeedRematchContract.lean});
\item co-reference-preserving canonical storage;
\item revision-based cache correctness;
\item incremental tabling soundness (\code{IncrementalTableSemantics.lean});
\item runtime contracts: canonicalization, rollback, effect classification (\code{CeTTaRuntimeContracts.lean});
\item theorem-to-runtime seam mapping.
\end{enumerate}

The seeded-rematch result is the key new contribution: CeTTa's imported backend is binding-correct because candidate narrowing (from PathMap trie or exact match) is followed by seeded native rematch (\code{space\_subst\_match\_with\_seed}).  The theorem \code{imported\_seedRematch\_binding\_parity} proves that when imported candidates are sound, the two-phase query (candidates then rematch) produces the same result set as direct matching.  Loop rejection (\code{seedRematchSafe\_no\_loop}) and seed extension (\code{seedRematch\_extends\_seed}) complete the safety story.

An implementation checklist (\code{imported-fast-path-checklist.md}) maps
each proved precondition to a specific CeTTa runtime check, identifies the
remaining blockers for enabling direct bridge bindings, and recommends a
phased enablement path (ground-only, single-variable, full V2).

As of April 2026, ground-only is enabled in the runtime (tracked by
\code{imported-v2-ground-skip}).  The next rung, single-variable queries,
is fully specified in Lean (\code{SingleVarSafe} in \code{ImportedRowContract.lean})
with conformance and negative tests, but remains gated pending bridge value
decode verification.  The Lean preconditions
(\code{singleVarSafe\_implies\_rematchFreeSafe}, \code{ground\_value\_no\_loop\_risk},
\code{singleBinding\_merge\_decidable}) document exactly what must hold before
a MORK implementer flips that switch.

By contrast, the dedicated \PathMap{} formalization paper should keep the
trie-specific material foregrounded:
\begin{enumerate}
\item algebraic hierarchy;
\item zipper structure;
\item concrete FTrie selector;
\item support/count and zero-crossing results insofar as they clarify trie
  usage.
\end{enumerate}

Similarly, the \CeTTa{} paper should use the results here to justify
runtime architecture choices, rather than reproving the whole formal
stack.

\section{End-to-End Pipeline}

The full formal pipeline is now connected across three Lean files:

\begin{center}
\begin{tabular}{@{}lll@{}}
\toprule
\textbf{Layer} & \textbf{File} & \textbf{Key Theorem} \\
\midrule
Bag semantics & \texttt{BagSupportBridge.lean} & \texttt{support\_bind} \\
Trie storage  & \texttt{HEBridge.lean} & \texttt{same\_support\_same\_trie} \\
Selection     & \texttt{CandidateArchitecture.lean} & \texttt{concrete\_twoPhase\_eq\_direct} \\
Matcher instantiation & \texttt{PathMapMatcherInstance.lean} & \texttt{he\_twoPhase\_eq\_direct} \\
Cache policy  & \texttt{CandidateArchitecture.lean} & \texttt{cache\_policy\_stable/growing/shrinking} \\
Variant reuse & \texttt{VariantQueryCorrectness.lean} & \texttt{variant\_queries\_same\_rhs} \\
Binding reuse & \texttt{BindingCacheSoundness.lean} & \texttt{variant\_cache\_binding\_shape} \\
Contract      & \texttt{CandidateArchitecture.lean} & \texttt{BackendContract}, \texttt{backendContract\_of\_faithful} \\
Pipeline      & \texttt{HEInterface.lean} & \texttt{end\_to\_end\_pipeline} \\
\bottomrule
\end{tabular}
\end{center}

The \texttt{end\_to\_end\_pipeline} theorem bundles all three properties in
one statement: given a faithful trie, (1)~query correctness holds,
(2)~no phantom candidates exist, and (3)~duplicate adds preserve support.

The \texttt{BackendContract} structure serves as the single proof obligation
for any \PathMap{}-backed space: provide \texttt{trieFaithful} (every
support atom has its encoding in the trie), and the full contract follows
automatically via \texttt{backendContract\_of\_faithful}.

The cache invalidation interface classifies mutations into three kinds:
\begin{itemize}
\item \texttt{supportStable}: no cache invalidation needed (duplicate adds,
  high-count removes);
\item \texttt{supportGrowing}: invalidate queries matching the new atom;
\item \texttt{supportShrinking}: invalidate queries matching the removed atom.
\end{itemize}

\noindent
Each maps to one branch in \texttt{table\_store.c}.

\section{Current Status and Next Proof Obligations}

The current development already supports the following claims:
\begin{enumerate}
\item support-level backend agreement is the correct first target;
\item \PathMap{} candidate selection plus native matching is a correct
  architecture;
\item duplicate adds do not need to mutate trie support;
\item injective canonicalization is the right shape for binder hygiene;
\item alpha-equivalent query keys preserve right-hand-side answers and
  result shape under cache reuse;
\item revision counters justify lazy cache invalidation.
\end{enumerate}

The most urgent next formal obligations are:
\begin{enumerate}
\item storage/materialization preserves query answers, not just variable
  injectivity;
\item snapshot/clone operations preserve backend faithfulness;
\item revision and snapshot semantics are connected explicitly enough to
  support staged or open table entries.
\end{enumerate}

These are the theorems most likely to help with the remaining nasty
binder bugs in the current \CeTTa{} runtime work.

A companion document (\texttt{theorem-runtime-matrix.md}) maps each proved
theorem to its runtime seam (\texttt{space.c}, \texttt{table\_store.c},
\texttt{space\_match\_backend.c}, \texttt{term\_universe.c}) and its
corresponding regression test target.  This matrix is the integration
contract between the formal development and the \CeTTa{} runtime.

\end{document}
